\documentclass[11pt]{article}

\usepackage[margin=1.1in]{geometry}
\usepackage{amsmath}
\usepackage{booktabs}
\usepackage{microtype}
\usepackage[round,authoryear]{natbib}
\usepackage[colorlinks=true,allcolors=blue]{hyperref}

\makeatletter
\renewcommand\paragraph{\@startsection{paragraph}{4}{\z@}%
  {3.25ex \@plus1ex \@minus.2ex}{-1em}%
  {\normalfont\normalsize\itshape}}
\makeatother

\title{\textbf{Beating a deployed rating system}\\[0.4em]
\large Partial Pooling Against Per-Category Glicko-2 on Lichess}
\author{Peter Cotton}
\date{\today}

\begin{document}
\maketitle

\begin{abstract}
\noindent Lichess maintains a separate Glicko-2 rating for each time
control, which is stratification: condition on the category, estimate
within it, and accept the sparse-cell variance that follows. We compare
it with a partially pooled alternative on 59{,}399 decisive games from
one month of the site's open database, in which a player's ability is a
level estimated from all of their games plus shrunk bullet and blitz
offsets. The pooled rating is better by $0.0137$ to $0.0182$ nats of
held-out log loss across three independent splits, every interval
excluding zero, and roughly four fifths of that margin is the
estimator, a batch fit against an online update, with one fifth the
factor structure itself. Stratification helps Glicko-2 and hurts our
scalar, because its benefit is capped by what pooling already captures
while its cost is not. Opening family is unidentifiable because players
choose openings, opponent-strength curvature is real but almost
entirely shared, and modern Lichess no longer poses the question, with
$0.3\%$ of players mixing time controls against $52\%$ in 2013.
\end{abstract}

\section{Introduction}

A rating is conventionally a number. Elo \citep{elo1978},
Glicko \citep{glickman1999}, TrueSkill \citep{herbrich2007} and their
descendants all summarise a competitor by a scalar, optionally with an
uncertainty attached. When a rating needs to reflect that ability
varies by context, say a player strong at speed chess and ordinary at
slow, the standard remedy is to stratify: keep one rating per
category. Lichess does exactly this, maintaining independent Glicko-2
ratings for bullet, blitz and classical.

Stratification is a choice with a price. Each category's rating is
estimated only from that category's games, so a player with many blitz
games and few bullet games gets a well-estimated blitz rating and a
noisy bullet one, and nothing the blitz games reveal about their
general strength is allowed to inform the bullet number. The
alternative is partial pooling: estimate the level from everything and
let only the \emph{deviations} be category-specific, shrunk toward
zero. Ability becomes a vector $\mu_i(x) = d_i \cdot x$ with
$x = [1, \mathbf{1}\{\text{bullet}\}, \mathbf{1}\{\text{blitz}\}]$,
and a per-feature penalty leaves the level free while shrinking the
offsets.

The question is which is better, measured against the deployed system
rather than a textbook baseline, and what the difference is made of.

\section{Data and protocol}

One month of the Lichess open database (2013-01, CC0), reduced to
headers. We keep decisive games between players with at least 100 games in the
month, 59{,}399 games among 639 players; draws are dropped. Each game is labelled with its
time control.

Games are split 50/25/25 into fitting, validation and scored sets, and
every arm sees the same games. Glicko-2's $\tau$ is tuned on the
validation split exactly as our ridge is; neither side is given a free
parameter the other lacks. Held-out mean log loss is reported on the
scored quarter only, and intervals come from a paired bootstrap over
20{,}000 resamples. Three independent splits are reported throughout,
because a single split on this kind of data is not evidence.

Four arms are compared. \emph{Glicko-2 pooled} ignores time control.
\emph{Glicko-2 per time control} is Lichess's deployed architecture.
\emph{Our scalar} is a batch Thurstonian maximum-a-posteriori fit with
one parameter per player, and exists to separate the estimator from the
structure. \emph{Our factor} adds the shrunk offsets.

\section{Result}

\begin{center}
\begin{tabular}{lccc}
\toprule
arm & seed 0 & seed 11 & seed 22 \\
\midrule
Glicko-2 pooled & 0.6323 & 0.6304 & 0.6365 \\
Glicko-2 per time control \emph{(Lichess)} & 0.6234 & 0.6241 & 0.6294 \\
our scalar & 0.6117 & 0.6131 & 0.6145 \\
\textbf{our factor} & \textbf{0.6073} & \textbf{0.6104} & \textbf{0.6112} \\
\bottomrule
\end{tabular}
\end{center}

\noindent The pooled factor rating beats the deployed architecture on
every split. Decomposed:

\begin{center}
\begin{tabular}{lccc}
\toprule
comparison & seed 0 & seed 11 & seed 22 \\
\midrule
\textbf{factor vs Lichess} & $-0.0161$ & $-0.0137$ & $-0.0182$ \\
estimator (scalar vs Glicko-2 pooled) & $-0.0206$ & $-0.0173$ & $-0.0219$ \\
structure (factor vs our scalar) & $-0.0044$ & $-0.0027$ & $-0.0033$ \\
stratification's value to Glicko-2 & $+0.0089$ & $+0.0063$ & $+0.0071$ \\
\bottomrule
\end{tabular}
\end{center}

\noindent Every interval excludes zero on every split, with
$P(\text{better}) \ge 0.995$.

The headline should not be quoted without the decomposition. Roughly
four fifths of the margin is the estimator, a batch fit that sees the
training games at once against an online update that sees each game
once, and one fifth is the factor structure. The structure's own
contribution, $0.003$ to $0.004$ with intervals excluding zero on all
three splits, is the part that answers the question posed. The estimator column is a real advantage but
a different claim, and it is partly a privilege: a production system
must process games once, and ours need not.

\section{Why stratification helps Glicko-2 and not the batch estimator}

Two facts in the table look contradictory and are not. Stratifying
Glicko-2 \emph{gains} $0.0063$--$0.0089$. Stratifying our scalar
\emph{loses} $0.0006$. Partial pooling of the same information gains
$0.0027$--$0.0044$.

Stratification's payoff runs inversely to the quality of the base
estimator. Its benefit is the category information; its cost is
sparse-cell variance. When the base estimator is weak, the benefit
dominates. When the base estimator is strong, the pooled level already
captures most of what the category would have told you, and only the
variance cost remains. Partial pooling extracts the same signal without
paying that cost, which is why it wins in both regimes and is the
unconditional recommendation. A practitioner reporting that
stratification helped them has reported something about their base
estimator, not found a counterexample.

\section{The covariates that did not work}

The time-control result is not evidence that any covariate would do,
and the failures are more informative than the success.

\paragraph{Opening family fails, and the mechanism generalises.} A fair
fit, with the offset ridge swept to the scalar limit, collapses exactly
onto the scalar. The axis is nonetheless real: among 790 players with at
least 20 games in each family, the style differential carries a true
standard deviation of about 4.7 win-rate points, roughly four sigma
above binomial noise. It is \emph{unidentifiable, not absent}. Players
choose their openings, so the within-player contrast that
identification requires barely exists: the tactical share is bimodal
with deciles $0.00/0.10/0.86/0.98/1.00$ and a standard deviation of
$0.415$, against $0.058$ under a no-choice null. A factor rating needs
a covariate that varies within a player for reasons other than the
player's preference.

\paragraph{Opponent strength is real but shared.} The signed
opponent-strength gap passes the exogeneity requirement by
construction, since pairing hands every player a spread of opponents.
Fitted against a deliberately fair baseline, a scalar plus one global
curvature coefficient rather than a bare scalar, the per-player offsets are identifiable, with a tuned ridge interior at
$10.0$ and surviving offset standard deviation $0.1316$. But the gain
is $-0.00056$ with interval $[-0.00117, +0.00005]$, spanning zero, and
eight times smaller than time control. The global coefficient is
$-0.584$: one shared parameter absorbs nearly all of the
opposition-strength signal, leaving only per-player deviations for the
factor arm. Against a bare scalar this axis would have looked strong,
and that comparison would have been a straw man.

So of three identifiable covariates, level, time control and opponent
strength, the factor treatment pays on one.

\section{Robustness}

\paragraph{Colour is not the confound.} Our arms model the white
advantage and standard Glicko-2 does not, so the estimator column
invites the suspicion that it is colour bookkeeping. It is not: colour
is worth $-0.0016$ of the $-0.0206$ gap, and a colour-blind fit still
beats colour-blind Glicko-2 by $-0.0190$ like for like.

\paragraph{The result is not an artifact of the dense subpopulation, but
the mechanism's own prediction failed.} Rerunning at inclusion
thresholds of $100/50/25/10$ games moves observations per parameter
from $30.9$ to $15.3$ and the margin stays between $-0.016$ and
$-0.019$ throughout, every interval excluding zero. The robustness
question is answered. The registered prediction, however, was that
pooling should pay \emph{most} where data is scarcest and the margin
should widen. It does not: it is flat within noise and not monotone.
The factor advantage is roughly constant in population density rather than concentrated among data-poor players,
which leaves the pooling account unrefuted but unconfirmed by the one
test that could have confirmed it.

\paragraph{Our arm is not winning by being overconfident.} Glicko-2
propagates rating deviation into every prediction; our arm prices a
point estimate with fixed noise, so only one side carries its own
uncertainty. Giving ours a validation-tuned temperature returns
$s^\star = 1.1$ and improves held-out loss by $0.0001$, with an
interval spanning zero; the margin moves from $-0.0161$ to $-0.0162$.
The arm was already calibrated and the asymmetry costs nothing. For
completeness, pricing instead with a per-player predictive variance
taken from the diagonal of the inverse penalised Hessian is
significantly \emph{worse}, by $+0.0015$.

\section{Limitations}

\paragraph{Batch against online.} Our estimator revisits the training set;
Glicko-2 processes each game once because a deployed system must. Part
of the estimator column is that privilege, and the deployment-realistic
comparison, online against online, is not run here.

\paragraph{Draws are dropped.} They are 3.3\% of games in this pool and far more at
longer time controls and higher strength. A three-way outcome model
is the right treatment.

\paragraph{One month of 2013.}
The obvious next step is a modern month, and it cannot currently be
taken. In 2024-01 the time-control covariate is unidentifiable: only
$0.3\%$ of players have at least 20 games in two or more time controls,
against $52\%$ in 2013, at every inclusion threshold we tried. Contrast
falls as activity rises, since the most active players are the most
specialised, 93.1\% bullet, so more data makes identification worse
rather than better, and restricting to the players who do mix controls
fails from the other side, leaving 904 players sharing 2{,}839 games
at roughly one observation per parameter. This is not evidence against
the result; no arm was fitted. It means the result holds on a
population where players mixed time controls, and that population no
longer exists at scale on this site.

\section{Conclusion}

A partially pooled vector rating beats a rating system in production,
on its own data, by $0.0137$ to $0.0182$ nats. Most of that is the
estimator and about a fifth is the pooling, and both parts are worth
stating separately. The wider claim the data supports is narrower than
``ratings should be vectors'': it is that \emph{conditioning} is the
expensive way to use category information and \emph{pooling} is the
cheap one, that the gap between them is largest for whoever has the
weaker base estimator, and that the covariate must vary within a player
for reasons the player does not choose. Opening choice fails that
condition, opponent strength satisfies it while carrying almost no
per-player signal, and time control satisfied it in 2013 and no longer
does.

\begin{thebibliography}{9}
\bibitem[Elo(1978)]{elo1978} A.~E. Elo. \emph{The Rating of Chessplayers,
Past and Present}. Arco, New York, 1978.
\bibitem[Glickman(1999)]{glickman1999} M.~E. Glickman. Parameter
estimation in large dynamic paired comparison experiments.
\emph{Journal of the Royal Statistical Society C} 48(3) (1999) 377--394.
\bibitem[Glickman(2012)]{glickman2012} M.~E. Glickman. Example of the
Glicko-2 system. Boston University technical note, 2012.
\bibitem[Herbrich et~al.(2007)Herbrich, Minka, and Graepel]{herbrich2007}
R.~Herbrich, T.~Minka, T.~Graepel. TrueSkill(TM): a Bayesian skill
rating system. \emph{Advances in Neural Information Processing Systems}
19 (2007) 569--576.
\bibitem[James and Stein(1961)]{james1961} W.~James, C.~Stein.
Estimation with quadratic loss. \emph{Proceedings of the Fourth Berkeley
Symposium on Mathematical Statistics and Probability} 1 (1961) 361--379.
\bibitem[Efron and Morris(1975)]{efron1975} B.~Efron, C.~Morris. Data
analysis using Stein's estimator and its generalizations. \emph{Journal
of the American Statistical Association} 70(350) (1975) 311--319.
\bibitem[Thurstone(1927)]{thurstone1927} L.~L. Thurstone. A law of
comparative judgment. \emph{Psychological Review} 34(4) (1927) 273--286.
\end{thebibliography}

\end{document}
